\documentclass[addpoints]{exam}
\pagestyle{headandfoot}

\usepackage[T1,T2A]{fontenc}
\usepackage[utf8]{inputenc}
\usepackage[russian]{babel}

\usepackage{amsmath, amsfonts}
\usepackage{verbatim}
\usepackage{graphicx}
\usepackage[usenames,dvipsnames]{color}
\usepackage{parskip}
\usepackage{tikz} % drawing schemes
\usepackage{hyperref} % references
\newcommand{\semester}{WS 2024/2025}
\runningheader{\students}{}{\semester}
\runningfooter{}{\thepage}{}
\headrule
\footrule

% ---------- Modify team name, students (matriculation number), exercise number here ----------
\newcommand{\teamname}{dl2023-team1}
\newcommand{\students}{Великанов Максим}
\newcommand{\assignmentnumber}{1}
% ---------- End Modify ----------

\title{mAP for beam search}
\author{Великанов Максим}
\date{\today}

\begin{document}
    \maketitle

    % ---------- Add Solution below here ----------

    The outline for computing the precision-recall curve and mAP is as follows:
    \begin{enumerate}
    	\item Make a list of all predictions. Confidence of a prediction is the log-probability of the beam under the model:\\
    	\verb|pred = [(episode_i, conf_i, traj_i)]; len(pred) = N_pred|\\
    	Make a list of true trajectories:\\
    	\verb|true = [(episode_i, 1.0, traj_i)]; len(true) = N_true|
    	\item Sort predictions by descending confidence
    	\item For each predicted trajectory:
    	\begin{enumerate}
    		\item calculate the L1 error with the true trajectory from the same episode
    		\item if \verb|error > error_thr| or another prediction was already found for this true trajectory, it is a false positive
    		\item otherwise, the prediction is a true positive
    	\end{enumerate}
    	\item Calculate \verb|precision_i| and \verb|recall_i| for the first \verb|i| predictions (which is equivalent to ignoring all predictions with \verb|conf_j < conf_i|). Precision is \verb|N_TP/N_pred|, recall is \verb|N_TP/N_true|.
    	\item Plot the (recall, precision) points and compute Area Under Curve (AUC)
    	\item Repeat the AUC calculation for several L1 error thresholds, the average AUC is the mean Average Precision
    \end{enumerate}
    
    % ---------- End of Document ----------
    
    \section{References}
    \begin{enumerate}
    	\item \href{https://jonathan-hui.medium.com/map-mean-average-precision-for-object-detection-45c121a31173}{Jonathan Hui}
    	\item \href{https://github.com/matin-ghorbani/Mean-Average-Precision-from-Scratch/blob/main/mAP.py}{code}
	    \item \href{https://learnopencv.com/mean-average-precision-map-object-detection-model-evaluation-metric/}{learnopencv - additional theory}
	    \item \href{https://pyimagesearch.com/2022/05/02/mean-average-precision-map-using-the-coco-evaluator/}{pyimagesearch - additional theory}
	\end{enumerate}

\end{document}
